\begin{pastexam}[label={px:gw3-097-ds-01}]{關務三等資料處理 97 年 第一題}
（申論；10 分）

說明演算法（Algorithm）的意義及詳述演算法必備的條件有那些？（10 分）
\end{pastexam}

\begin{pastsoln}
\textbf{意義。} 演算法是解決某一類問題的一組\textbf{有限、明確}的步驟：給它這類問題的任何一個合法
輸入，照著步驟一步一步做，一定會在有限步之後停下來，並產生正確的輸出。 它和程式的差別在於：演算法
是方法本身，可以用文字、虛擬碼或流程圖描述；程式是用某種程式語言把演算法寫出來、可以在電腦上執行
的東西（資結書第~1~章）。

\textbf{必備的五個條件}（Knuth 的說法，也是國考最常見的答法）：
\begin{enumerate}[nosep]
\item \textbf{輸入}（input）：有零個或多個由外界提供的輸入資料。
\item \textbf{輸出}（output）：至少產生一個輸出，也就是問題的答案。 不產生任何結果的步驟沒有意義。
\item \textbf{明確性}（definiteness）：每一個步驟的意思都清楚、沒有歧義，執行的人（或機器）不需要
自己猜。 例如「把 $x$ 加上一個適當的數」就不明確。
\item \textbf{有限性}（finiteness）：對任何合法輸入，執行有限步之後一定會結束。 例如「從 $1$ 開始
一直加一，直到找到最大的整數」永遠不會結束，不是演算法。
\item \textbf{有效性}（effectiveness）：每一個步驟都夠基本，原則上可以用紙筆在有限時間內做完。
例如「若哥德巴赫猜想成立則輸出 1」這種步驟，目前沒有人做得到，不算有效。
\end{enumerate}
另外通常也要求\textbf{正確性}：對每一個合法輸入，輸出都是正確的答案。

\textbf{例子：輾轉相除法求最大公因數}。 輸入兩個正整數 $a, b$（輸入）；若 $b = 0$ 輸出 $a$
（輸出），否則令 $(a, b) \gets (b, a \bmod b)$ 再重複（每一步都明確、而且是可以手算的基本運算）。
每做一次 $b$ 都嚴格變小，又不會小於 $0$，所以有限步後一定停（有限性）。 五個條件都滿足。
\end{pastsoln}
